\chapter{Sheaf Cohomology and Global Galaxy Maps}
\label{ch:global-galaxy-maps}

This chapter introduces the cohomological machinery required to construct
\emph{global semantic galaxy maps} from the local embedding data developed in
Chapters~\ref{ch:coupling}--\ref{ch:embedding-flow}.
The key idea is that Polyxan semantic neighborhoods form a sheaf
$\mathcal{F}$ on the hypergraph $\mathcal{G}$; 
RSVP semantic structures form a sheaf $\mathscr{S}$ on the manifold 
$\mathcal{M}$; and embeddings
\[
  \iota : \mathcal{G} \to \mathcal{M}
\]
induce a pullback sheaf 
\(
  \iota^{-1}\mathscr{S}
\)
that must glue consistently with $\mathcal{F}$.

Global galaxy maps arise precisely as \emph{global sections}
of the sheaf morphism
\[
  \mathcal{F} \longrightarrow \iota^{-1}\mathscr{S},
\]
and their existence is governed by cohomology, 
obstruction classes, and derived-gluing conditions.

This chapter provides:

\begin{itemize}
  \item a formal definition of galaxy maps as sheaf morphisms,
  \item Čech cohomology on $\mathcal{G}$ and $\mathcal{M}$,
  \item obstruction theory for gluing local semantic data,
  \item hypercover refinements for galaxy formation,
  \item homotopical and derived notions of semantic consistency,
  \item and stability theorems linking cohomology to embedding flows.
\end{itemize}

% ============================================================
\section{Local Sheaves: Polyxan and RSVP}
% ============================================================

Recall the two main semantic sheaves:

\begin{enumerate}
  \item The \emph{Polyxan sheaf} 
  \[
    \mathcal{F}: \mathrm{Open}(\mathcal{G})^{op} \to \mathbf{Cat}
  \]
  assigns to each sub-hypergraph $U\subseteq\mathcal{G}$ 
  a category of content atoms,
  with restriction functors given by deletion or projection of atoms.

  \item The \emph{RSVP semantic sheaf}
  \[
    \mathscr{S}: \mathrm{Open}(\mathcal{M})^{op} \to \mathbf{Cat}
  \]
  assigns to an open semantic region 
  a category of interpretations, field configurations,
  or semantic fibers (Chapter~\ref{ch:semantic-manifold}).
\end{enumerate}

Given an embedding $\iota:\mathcal{G}\to\mathcal{M}$,
we obtain the \emph{pullback sheaf}
\[
  \iota^{-1}\mathscr{S}: \mathrm{Open}(\mathcal{G})^{op}\to \mathbf{Cat}
\]
defined by
\[
  (\iota^{-1}\mathscr{S})(U)
  := \mathscr{S}(\iota(U)).
\]

A global semantic galaxy is precisely a compatible morphism between these sheaves.

% ============================================================
\section{Galaxy Maps as Sheaf Morphisms}
% ============================================================

\begin{definition}[Semantic Galaxy Map]
A \emph{semantic galaxy map} is a sheaf morphism
\[
  \Gamma : \mathcal{F} \longrightarrow \iota^{-1}\mathscr{S}
\]
consisting of natural transformations
\(
  \Gamma_U : \mathcal{F}(U) \to \mathscr{S}(\iota(U))
\)
compatible with all restriction maps.
\end{definition}

Intuitively:

\begin{itemize}
  \item $\mathcal{F}(U)$ encodes local Polyxan semantic data,
  \item $\mathscr{S}(\iota(U))$ encodes RSVP interpretations at embedded positions,
  \item $\Gamma_U$ assigns each content atom a semantic meaning at its geometric location.
\end{itemize}

Gluing all $\Gamma_U$ into a single global galaxy map demands nontrivial
cohomological conditions.

% ============================================================
\section{Čech Cohomology of the Polyxan Cover}
% ============================================================

Let $\mathfrak{U}=\{U_i\}$ be a good open cover of $\mathcal{G}$ by hypergraph neighborhoods.

Define the Čech cochains valued in $\iota^{-1}\mathscr{S}$:
\[
  \check{C}^k(\mathfrak{U},\iota^{-1}\mathscr{S})
  =
  \prod_{i_0,\dots,i_k}
    \mathscr{S}(\iota(U_{i_0}\cap\cdots\cap U_{i_k})).
\]

The Čech coboundary
\[
  \delta : \check{C}^k \to \check{C}^{k+1}
\]
encodes failure of compatibility on overlaps.

A global galaxy map exists iff the compatibility cocycle vanishes:

\begin{theorem}[Galaxy Map Existence Criterion]
A global section
\(
  \Gamma : \mathcal{F} \to \iota^{-1}\mathscr{S}
\)
exists iff the Čech 1-cocycle
\[
  \{\Gamma_{U_i}|_{U_i\cap U_j}
    - \Gamma_{U_j}|_{U_i\cap U_j}\}_{i,j}
\]
represents the zero class in
\[
  \check{H}^1(\mathfrak{U},\iota^{-1}\mathscr{S}).
\]
\end{theorem}

Thus:
\[
  \check{H}^1 = 0
  \quad\Longleftrightarrow\quad
  \text{global semantic galaxy exists and is unique}.
\]

This is the sheaf-theoretic analogue of solving constraint consistency in
knowledge graphs.

% ============================================================
\section{Obstruction Theory and Higher Cohomology}
% ============================================================

If $\check{H}^1$ is nonzero, the obstruction to global existence propagates.
Specifically, incompatibilities in triple overlaps yield a 2-cocycle:
\[
  \omega_{ijk}
  \in 
  \check{H}^2(\mathfrak{U},\iota^{-1}\mathscr{S}).
\]

\begin{definition}[Galaxy Obstruction Class]
The obstruction class of an embedding is
\[
  \mathfrak{o}(\iota)
  :=
  [\omega]
  \in \check{H}^2(\mathfrak{U},\iota^{-1}\mathscr{S}).
\]
\end{definition}

\begin{theorem}[Obstruction Criterion]
A global galaxy map exists 
\(
\iff
\mathfrak{o}(\iota)=0.
\)
\end{theorem}

This mirrors classical obstruction theory for principal bundles,
but now on a hypergraph–manifold correspondence.

% ============================================================
\section{Hypercover Refinement and Homotopical Consistency}
% ============================================================

In practice $\mathcal{G}$ may not admit a good Čech cover.
Thus we refine to \emph{hypercovers} in the sense of homotopy theory.

Let $\mathcal{U}_\bullet$ be a semisimplicial hypercover of $\mathcal{G}$.
Then derived sheaf cohomology is computed by
\[
  \mathbb{H}^k(\mathcal{G},\iota^{-1}\mathscr{S})
  =
  H^k(\mathrm{Tot}(\mathscr{S}(\iota(\mathcal{U}_\bullet)))).
\]

\begin{theorem}[Homotopical Gluing Theorem]
If $\mathbb{H}^1(\mathcal{G},\iota^{-1}\mathscr{S})=0$,  
then every partially defined galaxy map extends uniquely (up to equivalence) to a global galaxy.
\end{theorem}

This derived formulation is crucial for galaxy evolutions 
where embeddings distort the topology of $\iota(U)$ over time.

% ============================================================
\section{Relation to Embedding Flows (Chapter 9)}
% ============================================================

The embedding PDE of Chapter~\ref{ch:embedding-flow} 
pushes the Polyxan structure through RSVP potentials.
This modifies the pullback sheaf
\(
  \iota^{-1}\mathscr{S}
\),
and therefore changes the cohomology.

The dynamical system thus satisfies the following principle:

\begin{theorem}[Cohomological Attractor Principle]
Embedding flow $\dot\iota=-\delta\mathcal{F}/\delta\iota$  
tends to reduce 
\(
  \|\mathfrak{o}(\iota)\|
\)
monotonically (in the sense of derived obstruction norms).
\end{theorem}

\begin{proof}[Sketch]
The sheaf inconsistency term $\mathcal{E}_{\mathrm{sheaf}}$ in $\mathcal{F}$
acts precisely as the $L^2$ norm of the Čech obstruction cocycles.  
Steepest descent reduces these norms; thus cohomological defects decay.
\end{proof}

Embeddings are therefore \emph{cohomology-driven}:
semantic galaxies drift toward regions of the manifold
where global semantic gluing is possible.

% ============================================================
\section{Galaxy Maps as Fixed Points of Derived Functors}
% ============================================================

Let
\[
  \mathbf{R}\Gamma(\iota^{-1}\mathscr{S})
\]
denote the derived global-section functor.
Galaxy maps correspond to objects of
\[
  H^0(\mathbf{R}\Gamma(\iota^{-1}\mathscr{S})).
\]

\begin{theorem}[Fixed-Point Characterization]
A semantic galaxy map is a fixed point of the derived semantic functor:
\[
  \Gamma
  =
  \mathbf{R}\Gamma(\iota^{-1}\mathscr{S})(\Gamma).
\]
\end{theorem}

This gives a powerful categorical viewpoint:
galaxy maps are solutions of a homotopy-invariant fixed-point equation.

% ============================================================
\section{Large-Scale Structure and Semantic Cohomology}
% ============================================================

Cohomology determines the large-scale semantic geometry:

\begin{enumerate}
  \item $\check{H}^0$ classifies global semantic constants,  
  \item $\check{H}^1$ classifies gluing failures,  
  \item $\check{H}^2$ captures semantic curvature defects,  
  \item higher groups classify increasingly abstract holonomies  
        (ethical loops, narrative attractors, perceptual bundles).
\end{enumerate}

A semantic galaxy is “smooth’’ when higher cohomology vanishes.

We can therefore define:

\begin{definition}[Cohomologically Smooth Galaxy]
A semantic galaxy is smooth if  
\[
  \mathbb{H}^k(\mathcal{G},\iota^{-1}\mathscr{S})=0
  \quad\forall k\ge 1.
\]
\end{definition}

Smooth galaxies correspond to perfectly integrated, 
globally consistent semantic structures.

% ============================================================
\section{Summary}
% ============================================================

This chapter established the sheaf-cohomological foundations of
\emph{semantic galaxy maps}:

\begin{itemize}
  \item definition of galaxy maps as sheaf morphisms 
        $\mathcal{F} \to \iota^{-1}\mathscr{S}$,
  \item Čech cohomology and gluing criteria,
  \item obstruction classes in $H^2$,
  \item hypercover refinement and derived cohomology,
  \item fixed-point characterization in derived categories,
  \item coupling to embedding flows via sheaf inconsistency terms,
  \item interpretation of large-scale semantic structure via cohomology.
\end{itemize}

These constructions provide the global, topological substrate on which
semantic galaxies evolve.  
Chapters~\ref{ch:simulation-architecture} and \ref{ch:nbody}
will use this cohomological machinery to build
real-time simulation engines for galaxy morphogenesis.
